% !TeX root = ../papers/nc-p3.tex
\section{Walsh階数14窓の添字保持同型曲線と唯一の周期1点}
\label{sec:walsh-rank-fourteen-orbits}
\begin{theorem}\label{thm:walsh-rank-fourteen-orbits}
定理\ref{thm:walsh-symmetric-windows}の正則な階数14窓の添字保持同型類は
$\mathbb A^1\setminus\{-1,1\}$ の点と一対一に対応する。
$a=1,b=0$ の表示で $V=H^2$、
$E_6=16(L^2-M^2)(B^2-A^2)(D^2-C^2)$ と置けば座標は
$z=E_6/(4V)$、$V=4(1-z^2)$ である。一位置移動は $z\mapsto-z$ となり、
$z=0$ が唯一の最小周期1点、それ以外の全点の最小周期は2である。
\end{theorem}
\begin{proof}
$P_{h,u}$ は $(i,k)$ 成分が $[i=k+h](-1)^{u\cdot k}$ である四次のPauli行列とする。
反対切断を $\sum\gamma_{h+4u}P_{h,u}\otimes P_{h,\beta u}$ と書く。
$\beta$ はラベル順に $(0,3,2,1)$ である。二つの核はともに
$I,P_{1,0}$ で張られる。実際、この二ラベルだけが消え、Pauli基底のトレース双対は
各行列自身のスカラー倍である。各核内には階数2の直線が $I\pm P_{1,0}$ の
二つだけあり、いずれも二次元の像と相手側の二次元空間を結ぶ同型をスカラー倍まで定める。
これで各二頂点に結び付いた二次元直和分解が定まる。低位ビットをFourier変換すると、
二直線の代表は二直和因子それぞれの単位行列となる。

その分解で各行列因子の対角ブロック同士にテンソルを射影する。
異なる低位符号の組で射影は零、同じ低位符号の組では核の同型を通して
$\operatorname{End}(k^2)^{\otimes2}$ の元
\[
 2\big(A X\otimes X+ Z\otimes Z+B XZ\otimes XZ\big),\qquad
 X=\begin{pmatrix}0&1\\1&0\end{pmatrix},\quad Z=\operatorname{diag}(1,-1)
\]
を得る。従って偶位置と奇位置の二分解の対応も窓から決まる。
核の同型のスカラーと二つの基底変更は、上のテンソルに共通倍率と独立な共役で作用する。
各側のトレース形式で随伴との積を取ると、トレース零空間上の固有値の射影比は
両側で $[A^2:1:B^2]$ の順不同の組となる。
$A^2+B^2=1$, $AB=-H/4\ne0$ なので、他の二つの和に等しい値はただ一つである。
これを1に規格化すると残り二値の積が $V/16$ となる。よって $V$ は指定基底を忘れた不変量である。
さらに固有値1は単純であり、対応する無跡行列直線は $kZ$ である。
二固有線を取れば、四つの全頂点で各二空間がさらに二直線へ自然に分かれる。
$A^2=B^2$ のときも $kZ$ の固有値は単純なままである。

$H^2=V$ を固定し、$R^2=H/2$, $Q^2=1+R^2$, $N^2=1-R^2$, $i^2=-1$ を取る。
前節の三行列表示で
\[
 r=R,\quad l=iR,\quad q=Q,\quad n=N,\quad m'=iR,\quad p=\pm R
\]
の二代表を使う。この二代表と $b=0$ を保つ表示変換だけで、固定した $V$ の全ての窓を覆う。
実際、一般の $V$ では前節の復元で128組あり、各代表の64変換は相異なる。
$P_2=2(a^2+b^2+(a+b)^2)$ と $E_6$ は192変換で不変
（定理\ref{thm:rank-fifteen-rational-quotient}の斉次恒等式）なので、再規格化 $a'=\pm1$
も $E_6$ を変えない。二代表の $E_6$ は符号の異なる $4VQN$ である。
この被覆は特殊値でも成立する。固定した $R,Q,N$ に対する平方根の全組は、
$H$ を反転する場合の $R\mapsto iR$, $(Q,N)\mapsto(N,Q)$ も含めて
$R,Q,N$ の符号・四乗根倍で列挙される。一般値での被覆はこれらの線型式の恒等式なので
全点に特殊化できる。$V=4$ では前節により一軌道である。

残るのは二代表間の指定なし同型と一位置移動である。上の四直線の標識で、
各基底を低位ビット（分解の符号）と高位ビット（$Z$ の固有直線）で順序付ける。
同型の置換部分は $(P,Q',P,Q')$ の形で、$P,Q'$ は低位符号を同時に反転するかともに保ち、
各低位符号内の高位二直線を独立に交換できる。二つの核の階数2直線を単位行列にしたことから、
位置0,2の置換が一致し、位置1,3でも一致する。従って候補は32個で尽きる。
巡回したテンソルでも核の転置と二因子の交換によって同じ四直線を得る。

この基底での成分は、共通倍率を除き $1,R,Q,N$ の一つの四乗根倍である。
具体的に文字を $i',j,k,t$、$d=(i'+k)_0$, $g=(i'+k)_1$ とすれば、
$i'+k=j+t$ の成分は
\[
 \frac12\sum_{h,b'=0,1}\gamma_{h+2g+4(d+2b')}
   (-1)^{(i'+j)_0h+b'k_1+(b'+d)t_1}
\]
となる（ラベルビットの加法は標数2）。$QN\ne0$ では台は固定である。
台の16列の重み行列を作り、各語の行に対応する四位置の基底変数を記録する。
その整数左核の基底に沿った比の単項式が全て1であることが、対角基底変更の必要十分条件である。
Smith標準形と上の成分式で得る個数を示す。
\[
\begin{array}{c|rrr}
 (+R)\text{ 表示との比較先}&\text{台を保つ}&\text{全比が恒等的に1}&
 \text{定数比だけで不一致}\\ \hline
 -R&8&0&8\\
 (+R)\text{ の一位置移動}&8&0&8\\
 (-R)\text{ の一位置移動}&8&4&4\\
 (+R)\text{ の二位置移動}&8&4&4
\end{array}
\]
\path{computations/verify_walsh_rank_fourteen_orbits.py} は各比を
$i^eR^jQ^kN^l$ として計算する。最後の列では $(j,k,l)=(0,0,0)$,
$e\not\equiv0\bmod4$ の等式が要求されるので、どの非零特殊値でも不一致である。
$Q=0$ または $N=0$ では対応する成分を除いても、符号を替えた代表の一位置移動への
恒等的な対角同値が残る。両代表はそこで同型なので周期1である。
$E_6^2=16V^2(1-V/4)$ から所定の座標を得る。各値を実際に平方根で表示でき、
二代表の区別と移動が上表で確定したので主張が従う。
\end{proof}
